\section{Introduction}\label{sec:introduction}

Let $b(d)=\max(1,\lceil\log_2d\rceil)$ and $L=d^2b(d)$.
Question (T) asks whether one deterministic multitape Turing machine, with a fixed
finite number of tapes and a fixed finite alphabet, transposes all such matrices
in $O(L)$ time. The input and output are written in row-major order, so the
output lists the original columns consecutively. The question is motivated in
part by the reductions between transposition and integer multiplication of
Harvey and van der Hoeven~\cite{HvdH2025}.

An obstruction to routing alone does not settle (T). A tape computation may
store nonlinear functions of many input bits and later combine several such
records. The first paper of this series isolates that distinction
\cite{FirstPaper2026}. The depth-reduction program of the fourth and eighth
papers cuts a dependency graph by revealing a small transcript, and then asks
for a coding bound on the remaining recovery relation
\cite{FourthPaper2026,EighthPaper2026}. The seventh paper proves half-rate
bounds under explicit elimination hypotheses~\cite{SeventhPaper2026}.
This paper instead constructs complete source certificates for particular
return geometries and proves that their information deficit exceeds the
cost of the revelations.

\subsection{The argument and its scope}

Fix one full computation geometry and a nonempty set $\mathcal C$ of inputs
correctly transposed on that geometry. Write
\[
 \Delta=L-\log_2|\mathcal C|.
\]
A transcript of at most $g$ bits partitions $\mathcal C$ into at most $2^g$
fibres. If, on every fibre, the complementary coordinates determine a set
$A$ of input coordinates, then
\begin{equation}\label{eq:intro-paid-deficit}
 \Delta\ge |A|-g.
\end{equation}
Our work is to supply $A$ and $g$ from the same actual geometry. Complete
source sets retain direct input reads, state and work-tape parents, and
unpaid parents entering an output vertex whose record has been revealed.
The certificate permits nonlinear computations and correlated correctness
sets. It never assigns independent information to different source paths.

The main geometric hypothesis concerns complete return families. A
$\mathrm{Q}_2$ core has two classes of pairwise disjoint full return spans;
the definition and the strict clipping convention are given in
\cref{sec:source-engine}. An exceptional family is one complete family,
not one work tape. In addition to this decomposition we retain fixed
runtime and whole-step capacity constants and a fixed bound on the number
of disjoint-span classes in the full family system.

\begin{enumerate}
\item \emph{A finite source certificate.} A single cyclic ordering of the raw
coordinates produces an acyclic recoverable set. The size of that set is
controlled by the circular gap areas of complete source banks, with explicit
movement, grid and output-stripe errors. Fractional covers and anchored
approximation improve its geometric term without changing the recovery bank.
\item \emph{Uniform restricted impossibility.} A $\mathrm{Q}_2$ core plus one
arbitrary exceptional family cannot have both $\Delta=o(L)$ and maximum
outside-family size $o(L)$. For $q$ fixed exceptions, every limiting sorted
hop profile has tail at least a positive multiple of $R^{1-q}$.
Consequently a common size envelope $H/i^p$ with $p>q$ is excluded.
\item \emph{Further sufficient conditions.} Economical fractional source
covers, or sufficiently shared complete raw support paid for once across
all outputs, give additional uniform correctness-deficit theorems.
The quantifiers keep every finite cut parameter fixed before dimension
increases. These are criteria on a native class, not assertions that an
arbitrary machine satisfies them.
\item \emph{Construction restrictions and limits.} We exhibit a quadratic
pointwise cover obstruction, prove a movement bound for a specified query
interface even after reordering, and show that isolating fixed groups of
$o(\log d)$ consecutive columns causally requires an almost complete helper.
Each result states the interface it constrains.
\end{enumerate}

For an oblivious transposer correct on every input, $\Delta=0$; a positive
uniform deficit therefore excludes its geometry. For an adaptive machine,
passing to one geometry may already discard much of the input entropy.
We do not infer a large common-geometry correctness set without paying
that loss. Nor do we prove that the stronger cover criteria exclude a
native class outside the union of all earlier exclusions.

\subsection{Relation to the earlier papers and other work}

The source-bank construction uses actual return families and the owner
ordering studied in the eighth paper~\cite{EighthPaper2026}, and the
dependency semantics of the fourth~\cite{FourthPaper2026}. We give the
additional packing, residual-cut and raw-source proofs here.
The fifth paper develops complete recovery relations and geometric
restrictions on their codes~\cite{DirectedRecovery2026}. Its distinction
between a complete relation and a convenient list of selected paths is
essential throughout the present argument. The local coding barriers of
the second paper~\cite{SecondPaper2026} concern different functional
interfaces; they do not supply the source covers assumed here.

The fixed-geometry deficit theorems are complementary to the depth
reductions of the eighth paper. A depth bound alone leaves a coding
question. Under our different geometric hypotheses we complete the
deficit calculation. This is not a general improvement of the seventh
paper's half-rate theorem or a proof of the unrestricted coding crux.

The movement example uses an established placement delivery array from
coded caching~\cite{YanChengTangChen2017}; we credit that code and prove
only the stated additional query-movement assertion. The causal-group
argument uses finite max-flow/min-cut~\cite{FordFulkerson1956} and the
first paper's native routing metric. Its additional content is a paid
extraction of partial routing from a locally autonomous computation on
correlated helper fibres. We make no claim that the elementary covering,
entropy or flow principles are new.

\subsection{Proof status and organization}

All 32 numbered results and 6 numbered definitions
are formalized in Lean. They reuse the eighth paper's family geometry,
shared incidence, interval-cover, interval-graph, finite-orbit,
shifted-grid, causal-recovery, circular-gap and stripe-counting lemmas,
the second paper's finite-entropy library, the first paper's native
counting library, and Mathlib's compactness, summability and bounded-limit results. Their Lean statements are
authoritative; the accompanying outlines are guides to those proofs.
Cited results state their provenance and formal status
where used. The index in \cref{sec:theorem-index} identifies every exact statement,
its current status and the verified declarations. The index itself is not a
verification certificate.

\Cref{sec:source-engine} establishes the model, source bank and finite
freezing bound. \Cref{sec:source-covers} supplies the fractional and approximate
cover theorem. \Cref{sec:uniform} derives the uniform profile consequences;
\cref{sec:shared-support} treats raw support shared across outputs.
\Cref{sec:limitations,sec:movement,sec:causal-groups} explain the limits and
additional construction restrictions. \Cref{sec:roadmap} identifies the
remaining native geometry and computation questions.
